% Stage 11 — 2.1 Function Spaces and Maxwell Operators (Batch A reconstruction)
\section{Function Spaces and Maxwell Operators}
\label{sec:ch21}

Finite horns, ground-plane slots, and compact-range scan surfaces all pose the same mathematical question before any pattern or synthesis algorithm runs: which phasor fields and impressed currents are physically admissible on a declared domain \(\Omega\), and on which function classes do the Maxwell maps \(\mathcal{S}_\omega\), \(\mathcal{L}_\omega\), and \(\mathcal{R}_\omega=\Pi_{\mathrm{rad}}\circ\mathcal{S}_\omega\) act without contradiction \cite{stratton1941,harrington2001}? A rim current on a horn, a MoM sheet on a patch, or a finite-element port that fixes \(\hat{\mathbf{n}}\times\E\) on a conductor each supplies boundary data and sources that must belong to spaces where curl, divergence, and tangential trace are well defined. Section~\ref{sec:ch201} named those maps at the level of \(\mathcal{J}_\Omega\) and \(\mathcal{F}_\Omega\); the present section supplies the Sobolev, trace, and dual structure that makes domains and ranges precise enough for Green representations, outgoing selection, and spectral analysis later in the chapter \cite{maxwell1873}.

The finite-energy floor and admissible-state class imported from Chapter~\ref{ch:01}---Eqs.~\eqref{eq:ch1.energy-class} and~\eqref{eq:ch1.admissible-space} together with the source--boundary tuple of Subsection~\ref{sec:ch01.1.2}---are not re-derived here \cite{stratton1941}. Every space declaration below must remain compatible with harmonic charge continuity Eq.~\eqref{eq:ch2.continuity-operator-domain} and with the weak form anchored in Subsection~\ref{sec:ch203} through Eqs.~\eqref{eq:ch2.maxwell-energy-form} and~\eqref{eq:ch2.maxwell-weak-form}. Dyadic Green kernels, Sommerfeld outgoing conditions, radiation integrals, and modal spectra are developed in Sections~\ref{sec:ch23}--\ref{sec:ch26}; attempting to invert \(\mathcal{L}_\omega\) through a kernel before trace classes and dual pairings are fixed ignores boundary rows and corrupts flux identities \cite{stratton1941}.

Subsection~\ref{sec:ch211} defines source, field, trace, and distributional classes for phasor Maxwell theory on Lipschitz domains. Subsection~\ref{sec:ch212} treats linear operators---domain, range, graph norm, and adjoint pairing Eq.~\eqref{eq:ch2.maxwell-adjoint-pairing}. Subsection~\ref{sec:ch213} develops compactness and rank-deficient observation maps that foreshadow finite-aperture limits and \(\ker\mathcal{R}_\omega\) in Section~\ref{sec:ch27} \cite{chu1948,devaney2004}. Subsection~\ref{sec:ch214} addresses completions, passivity, and Maxwell-closed screens under which truncated operators inherit coercivity on bounded domains while exterior problems retain outgoing subspaces rather than full self-adjointness \cite{harrington2001}.

Two distinct rank limitations must remain separate throughout. Instrument maps---finite apertures, probe coupling, truncated modal bases---are finite-rank on field space even when \(\mathcal{L}_\omega\) is injective on the physical branch. Equivalent source distributions in \(\ker\mathcal{R}_\omega\) share the same far-zone export while differing in reactive near fields; that export null space is analyzed in Section~\ref{sec:ch27}, but the compact embeddings used there are prepared in Subsection~\ref{sec:ch213}. Measurement compositions from Chapter~\ref{ch:01} therefore appear as rank-deficient maps on the right of Eq.~\eqref{eq:ch2.continuity-composition-law}, not as modifications of Maxwell's local laws \cite{stratton1941}. Piecewise-constant constitutive data enter through the same product Sobolev template, with interface traces from Eq.~\eqref{eq:ch1.interface-conditions} enforced on \(\Gamma=\partial\Omega\) and PEC or Leontovich closures modifying the boundary pairing in Subsection~\ref{sec:ch214} \cite{stratton1941}.
